% from ega2-1-fr.tex
\oldpage[II]{122}

\begin{corollary}[6.4.4]
\label{II.6.4.4-fr}
Sous les hypothèses de (\hyperref[II.6.4.3-fr]{6.4.3}), si on suppose en outre
$A$ normal, les $\sigma_i(u)$ (en particulier $\operatorname{Tr}u$ et
$\det u$) appartiennent à $A$.
\end{corollary}

\begin{proposition}[6.4.5]
\label{II.6.4.5-fr}
Soient $A$ un anneau intègre, $E$ un $A$-module de type fini, de rang $n$, $u$
un endomorphisme de $E$ tel que les $\sigma_i(u)$ appartiennent à $A$. Pour
que $u$ soit un automorphisme de $E$, il est nécessaire que $\det u$ soit
inversible dans $A$~; cette condition est suffisante lorsque $E$ est sans
torsion.
\end{proposition}

La condition est \emph{suffisante}, car l'hypothèse et les formules
(\hyperref[II.6.4.1.4-fr]{6.4.1.4}) et
(\hyperref[II.6.4.1.5-fr]{6.4.1.5}) (valables dans $E$, et non seulement dans
$E\otimes_AK$, puisque $E$ est sans torsion) prouvent que
$(\det u)^{-1}Q(u)$ est l'inverse de $u$.

La condition est \emph{nécessaire}, car si $u$ est inversible, il résulte de
(\hyperref[II.6.4.3-fr]{6.4.3}) que $\det(u^{-1})$ appartient à la clôture
intégrale $A'$ de $A$ dans son corps des fractions $K$ et est évidemment
inverse de $\det u$ dans $A'$. Notre assertion résulte alors du~:

\begin{lemma}[6.4.5.1]
\label{II.6.4.5.1-fr}
Soit $A$ un sous-anneau d'un anneau $A'$ tel que $A'$ soit entier sur $A$. Si
un élément $x\in A$ est inversible dans $A'$, il est aussi inversible dans
$A$.
\end{lemma}

Dans le cas contraire, $x$ appartiendrait à un idéal maximal $\mathfrak{m}$
de $A$ et il résulte du premier th. de Cohen-Seidenberg
(\cite[t.~I, p.~257, th.~3]{I-13}) qu'il y aurait un idéal maximal
$\mathfrak{m}'$ de $A'$ tel que $\mathfrak{m}=\mathfrak{m}'\cap A$~; on aurait
donc $x\in\mathfrak{m}'$, ce qui est absurde.

\begin{corollary}[6.4.6]
\label{II.6.4.6-fr}
Soient $A$ un anneau intègre et intégralement clos, $E$ un $A$-module de type
fini sans torsion, $u$ un endomorphisme de $E$. Pour que $u$ soit un
automorphisme de $E$, il faut et il suffit que $\det u$ soit inversible dans
$A$.
\end{corollary}

Cela résulte de (\hyperref[II.6.4.4-fr]{6.4.4}) et
(\hyperref[II.6.4.5-fr]{6.4.5}).

\begin{remark}[6.4.7]
\label{II.6.4.7-fr}
Nous aurons besoin plus loin d'une généralisation des résultats précédents.
Considérons un anneau noethérien \emph{réduit} $A$~; soient
$\mathfrak{p}_\alpha$ ($1\leq\alpha\leq r$) ses idéaux minimaux, $K_\alpha$
le corps des fractions de l'anneau intègre $A/\mathfrak{p}_\alpha$, $K$
l'anneau total des fractions de $A$, \emph{composé direct} des corps
$K_\alpha$. Soit $E$ un $A$-module de type fini, et \emph{supposons que}
$E\otimes_AK$ soit un $K$-module libre de rang $n$ (ce qui ici n'est plus
conséquence des autres hypothèses)~; il revient au même de dire que tous les
$K_\alpha$-espaces vectoriels $E\otimes_AK_\alpha=E_\alpha$ ont \emph{même
dimension $n$}~; si alors $u$ est un endomorphisme de $E$, on pose encore
$P(u,T)=P(u\otimes1,T)$ et $\sigma_j(u)=\sigma_j(u\otimes1)$, et en
particulier $\det u=\det(u\otimes1)$~; les $\sigma_j(u)$ sont donc des
éléments de $K$. Il est immédiat que $E\otimes_AK$ est somme directe des
$E_\alpha$ et que ces derniers sont stables par $u\otimes1$~; la restriction
de $u\otimes1$ à $E_\alpha$ n'est autre d'ailleurs que l'extension $u_\alpha$
de $u$ à ce $K_\alpha$-espace vectoriel~; on en conclut que $\sigma_j(u)$ est
l'élément de $K$ dont les composantes dans les $K_\alpha$ sont les
$\sigma_j(u_\alpha)$. Comme la fermeture intégrale de $A$ dans $K$ est
\emph{composée directe} des fermetures intégrales de $A$ dans les
$K_\alpha$, les $\sigma_j(u)$ sont \emph{entiers} sur $A$.

\begin{lemma}[6.4.7.1]
\label{II.6.4.7.1-fr}
La sous-$A$-algèbre de $K$ engendrée par tous les éléments $\sigma_j(u)$
($1\leq j\leq n$) lorsque $u$ parcourt $\operatorname{Hom}_A(E,E)$, est un
$A$-module de type fini.
\end{lemma}

Il suffit de démontrer que la sous-$A[T]$-algèbre de $K[T]$ engendrée par les
$P(u,T)$ est un $A[T]$-module de type fini, car si les $F_i(T)$
($1\leq i\leq m$) forment un système de générateurs de ce $A[T]$-module, les
coefficients des $F_i(T)$ sont entiers sur $A$, donc engendrent une
$A$-algèbre qui est un $A$-module de type fini
(\cite[t.~I, p.~255, th.~1]{I-13}). On peut donc

